\chapter*{Introduction to This Volume: How Carbon Came Alive}
\addcontentsline{toc}{chapter}{Introduction to This Volume}
\markboth{Introduction to This Volume}{}

At the end of Volume I, we know how electron structure arranges atoms into the periodic table, how atoms bond, and what determines a reaction's direction and speed. The same rules apply to salt, iron nails, gasoline, and seawater.

Volume II begins with a special element: carbon. Its four valence electrons let it build chains, rings, branches, and networks. All of organic chemistry, and all the molecular materials of life, rest on this four-handed atom.

This volume has three stages. Part V explores how carbon builds the molecular world, from petroleum to alcohol, aspirin to plastics. You will discover that organic reactions need not be memorized one by one: the behavior of functional groups and electrons tells you what molecules tend to do. Part VI enters biochemistry: how carbohydrates, lipids, proteins, and enzymes make up life; how ATP transfers energy like a currency; and how glycolysis, respiration, and photosynthesis link food, sunlight, and oxygen in a network. Part VII zooms out one level: how these molecules assemble into cells, chemical cities that maintain themselves, receive messages, cooperate, and sometimes rebel.

Keep two tools from Volume I ready: the energy account, in which reactions favor lower energy and greater dispersal, and the rate account, in which activation energy controls speed while catalysts change the route without changing the overall balance. You will find that life does not violate these accounts. It uses them with extraordinary ingenuity.
